\documentclass{article}
\usepackage{amsmath,amssymb,amsthm}
\usepackage{algorithm}
\usepackage{algpseudocode}
\usepackage{hyperref}
\usepackage{enumitem}
\newtheorem{theorem}{Theorem}
\newtheorem{lemma}{Lemma}
\newtheorem{assumption}{Assumption}
\newcommand{\EE}{\mathbb{E}}
\newcommand{\RR}{\mathbb{R}}
\newcommand{\norm}[1]{\left\lVert#1\right\rVert}
\newcommand{\prox}{\operatorname{prox}}
\begin{document}

\section*{Algorithm Name}
\textbf{Stochastic Variance‑Reduced Gradient (SVRG) for Non‑Convex Smooth Optimization}

\section*{Mathematical Formulation}
We consider the finite‑sum problem  
\[
\min_{x\in\RR^{d}} f(x)\;,\qquad 
f(x)=\frac{1}{n}\sum_{i=1}^{n} f_i(x),
\]
where each component $f_i$ is differentiable and $L$‑smooth (definition given below).  

SVRG proceeds in \emph{epochs}.  
At the beginning of epoch $s$ a \emph{snapshot} point $\tilde{x}^{\,s}$ is fixed and the \emph{full gradient}
\[
\tilde{\mu}^{\,s}:=\nabla f(\tilde{x}^{\,s})=\frac{1}{n}\sum_{i=1}^{n}\nabla f_i(\tilde{x}^{\,s})
\tag{1}
\]
is computed once.  

Inside epoch $s$ we perform $m$ inner updates indexed by $t=0,\dots,m-1$.
Let $x_{0}^{\,s}:=\tilde{x}^{\,s}$.  For each inner iteration we sample an index $i_t$ uniformly from $\{1,\dots,n\}$ and construct the \emph{variance‑reduced estimator}
\[
v_t^{\,s}:=\nabla f_{i_t}(x_{t}^{\,s})-\nabla f_{i_t}(\tilde{x}^{\,s})+\tilde{\mu}^{\,s}.
\tag{2}
\]
The iterate is updated by a standard SGD step with constant stepsize $\eta>0$:
\[
x_{t+1}^{\,s}=x_{t}^{\,s}-\eta\,v_t^{\,s}.
\tag{3}
\]
After $m$ inner steps we set the next snapshot $\tilde{x}^{\,s+1}=x_{m}^{\,s}$ and repeat.

\section*{Pseudocode}
\begin{algorithm}[H]
\caption{SVRG for Non‑Convex $L$‑smooth Objectives}
\begin{algorithmic}[1]
\State \textbf{Input:} stepsize $\eta>0$, inner loop length $m\in\mathbb{N}$, total epochs $S$.
\State Initialize $ \tilde{x}^{\,0}\in\RR^{d}$ arbitrarily.
\For{$s=0,\dots,S-1$}
    \State Compute full gradient $\displaystyle \tilde{\mu}^{\,s}\gets\frac{1}{n}\sum_{i=1}^{n}\nabla f_i(\tilde{x}^{\,s})$  \Comment{cost $O(n)$}
    \State $x_{0}^{\,s}\gets\tilde{x}^{\,s}$
    \For{$t=0,\dots,m-1$}
        \State Sample $i_t\sim\text{Uniform}\{1,\dots,n\}$
        \State $v_t^{\,s}\gets\nabla f_{i_t}(x_{t}^{\,s})-\nabla f_{i_t}(\tilde{x}^{\,s})+\tilde{\mu}^{\,s}$
        \State $x_{t+1}^{\,s}\gets x_{t}^{\,s}-\eta\,v_t^{\,s}$
    \EndFor
    \State $\tilde{x}^{\,s+1}\gets x_{m}^{\,s}$
\EndFor
\State \textbf{Output:} Choose $\hat{x}$ uniformly at random from $\{x_{t}^{\,s}\mid s=0,\dots,S-1,\,t=0,\dots,m-1\}$.
\end{algorithmic}
\end{algorithm}

\section*{Dry Run Example}
We minimize the one‑dimensional quadratic $f(w)=(w-3)^2$.
Here $n=1$, $f_1(w)= (w-3)^2$, $\nabla f_1(w)=2(w-3)$, and $L=2$.

\begin{enumerate}[label=\arabic*.]
\item \textbf{Initialization:} $w_0=0$, stepsize $\eta=0.1$, inner length $m=3$, epochs $S=1$.
\item \textbf{Full gradient at snapshot:}
\[
\tilde{\mu}^{\,0}= \nabla f_1(\tilde{w}^{\,0}) = 2(0-3) = -6.
\]
\item \textbf{Inner iterations (since $n=1$ we always sample $i_t=1$).}
\begin{itemize}
\item $t=0$: $v_0 = \nabla f_1(w_0)-\nabla f_1(\tilde{w}^{\,0})+\tilde{\mu}^{\,0}=(-6)-(-6)+(-6)=-6$\\
$w_1 = w_0 -\eta v_0 = 0-0.1(-6)=0.6$,  $f(w_1)=(0.6-3)^2=5.76$.
\item $t=1$: $\nabla f_1(w_1)=2(0.6-3)=-4.8$\\
$v_1 = (-4.8)-(-6)+(-6) = -4.8+6-6 = -4.8$\\
$w_2 = w_1-0.1(-4.8)=0.6+0.48=1.08$, $f(w_2)=(1.08-3)^2=3.6864$.
\item $t=2$: $\nabla f_1(w_2)=2(1.08-3)=-3.84$\\
$v_2 = (-3.84)-(-6)+(-6)= -3.84+6-6 = -3.84$\\
$w_3 = w_2-0.1(-3.84)=1.08+0.384=1.464$, $f(w_3)=(1.464-3)^2=2.3637$.
\end{itemize}
\item \textbf{Snapshot update:} $\tilde{w}^{\,1}=w_3=1.464$.
\end{enumerate}
The objective value decreased from $9$ at $w_0$ to $2.36$ after three inner steps, illustrating the variance‑reduced correction.

\section*{Assumptions}
\begin{assumption}[Smoothness]\label{ass:smooth}
Each component $f_i:\RR^{d}\to\RR$ is $L$‑smooth, i.e.
\[
\|\nabla f_i(x)-\nabla f_i(y)\|\le L\|x-y\|\qquad\forall x,y\in\RR^{d}.
\tag{A1}
\]
\end{assumption}
\begin{assumption}[Bounded Variance of Stochastic Gradients]\label{ass:var}
For any $x\in\RR^{d}$,
\[
\frac{1}{n}\sum_{i=1}^{n}\|\nabla f_i(x)-\nabla f(x)\|^{2}\le\sigma^{2}.
\tag{A2}
\]
\end{assumption}
\begin{assumption}[Lower Bounded Objective]\label{ass:lb}
There exists $f^{\star}>-\infty$ such that $f(x)\ge f^{\star}$ for all $x$.
\tag{A3}
\end{assumption}
All constants $L,\sigma$ are known, and the stepsize $\eta$ will be chosen as a function of $L,m$.

\section*{Main Convergence Theorem}
\begin{theorem}[Non‑Convex SVRG Convergence]\label{thm:svrg}
Assume \ref{ass:smooth}–\ref{ass:lb} and pick a constant stepsize
\[
\eta=\frac{1}{3L}\quad\text{and}\quad m=2n.
\]
Let $\{x_{t}^{\,s}\}$ be the iterates generated by SVRG and let $\hat{x}$ be the output (uniform random pick).  Then
\[
\EE\bigl[\|\nabla f(\hat{x})\|^{2}\bigr]
\;\le\;
\frac{3L\bigl(f(\tilde{x}^{\,0})-f^{\star}\bigr)}{T}
\;+\;
\frac{2\sigma^{2}}{n},
\tag{4}
\]
where $T=S\cdot m$ is the total number of stochastic gradient evaluations.
Consequently, to obtain $\EE[\|\nabla f(\hat{x})\|^{2}]\le\varepsilon$ it suffices to use
\[
T=O\!\left(\frac{L\Delta_f}{\varepsilon}+ \frac{n\sigma^{2}}{\varepsilon}\right)
\]
oracle calls, with $\Delta_f:=f(\tilde{x}^{\,0})-f^{\star}$.
\end{theorem}

\section*{Theoretical Properties}
\begin{itemize}
\item \textbf{Stability w.r.t. stepsize.} Condition $\eta\le 1/(3L)$ guarantees that the Lyapunov function used in the proof is non‑increasing; larger stepsizes may break the descent property.
\item \textbf{Inner‑loop length $m$.} Choosing $m=2n$ balances the cost of a full gradient ($O(n)$) with the variance reduction effect; any $m=\Theta(n)$ yields the same $O(1/T)$ rate.
\item \textbf{Comparison to plain SGD.} For non‑convex $L$‑smooth problems, SGD with stepsize $\eta=O(1/\sqrt{T})$ attains $\EE[\|\nabla f(\hat{x})\|^{2}]=O(1/\sqrt{T})$.  SVRG improves this to $O(1/T)$ while keeping the per‑iteration cost comparable.
\item \textbf{Effect of variance bound $\sigma^{2}$.} The term $2\sigma^{2}/n$ reflects the residual variance after variance reduction; it vanishes when $n\to\infty$ (online setting) or when all $f_i$ are identical.
\end{itemize}

\section*{Discussion}
The theorem shows that even without convexity, the variance‑reduced estimator enables a linear decay of the expected squared gradient norm.  The proof relies on the smoothness inequality, the unbiasedness of $v_t^{\,s}$, and a carefully crafted Lyapunov function that couples the distance to the snapshot with the objective value.  The constants in (4) are explicit, which is useful for practical stepsize tuning.

\section*{Preliminaries (Definitions and Lemmas)}
\begin{definition}[Lyapunov function]\label{def:lyap}
For epoch $s$ and inner step $t$ define
\[
\Phi_{t}^{\,s}:= \EE\!\left[ f(x_{t}^{\,s}) + \frac{c}{2}\|x_{t}^{\,s}-\tilde{x}^{\,s}\|^{2}\right],
\tag{5}
\]
where $c>0$ will be specified later.
\end{definition}
\begin{lemma}[Smoothness inequality]\label{lem:smooth}
If $f$ is $L$‑smooth then for any $x,y$,
\[
f(y)\le f(x)+\langle\nabla f(x),y-x\rangle+\frac{L}{2}\|y-x\|^{2}.
\tag{6}
\]
\end{lemma}
\begin{lemma}[Unbiasedness of the estimator]\label{lem:unbiased}
For any $t,s$,
\[
\EE_{i_t}[v_t^{\,s}\mid x_{t}^{\,s},\tilde{x}^{\,s}]
= \nabla f(x_{t}^{\,s}).
\tag{7}
\]
\end{lemma}
\begin{proof}
\[
\EE_{i_t}[v_t^{\,s}]
=\frac{1}{n}\sum_{i=1}^{n}\bigl(\nabla f_i(x_{t}^{\,s})-\nabla f_i(\tilde{x}^{\,s})\bigr)+\tilde{\mu}^{\,s}
= \nabla f(x_{t}^{\,s})-\nabla f(\tilde{x}^{\,s})+\nabla f(\tilde{x}^{\,s})
= \nabla f(x_{t}^{\,s}).
\tag{8}
\]
\end{proof}

\section*{Main Proof (Step‑by‑Step Derivation)}
We prove Theorem~\ref{thm:svrg}.  Throughout the proof expectations are taken over all random choices up to the current iteration.

\subsection*{Step 1: One‑step descent of the Lyapunov function}
From Lemma~\ref{lem:smooth} with $x=x_{t}^{\,s}$ and $y=x_{t+1}^{\,s}=x_{t}^{\,s}-\eta v_t^{\,s}$,
\[
f(x_{t+1}^{\,s})\le f(x_{t}^{\,s})-\eta\langle\nabla f(x_{t}^{\,s}),v_t^{\,s}\rangle
+\frac{L\eta^{2}}{2}\|v_t^{\,s}\|^{2}.
\tag{9}
\]
\textbf{[Step 1]} Taking expectation conditioned on $x_{t}^{\,s}$ and using Lemma~\ref{lem:unbiased},
\[
\EE\!\bigl[f(x_{t+1}^{\,s})\mid x_{t}^{\,s}\bigr]
\le f(x_{t}^{\,s})-\eta\|\nabla f(x_{t}^{\,s})\|^{2}
+\frac{L\eta^{2}}{2}\EE\!\bigl[\|v_t^{\,s}\|^{2}\mid x_{t}^{\,s}\bigr].
\tag{10}
\]

\subsection*{Step 2: Bounding the second moment of $v_t^{\,s}$}
Write $v_t^{\,s}= \nabla f_{i_t}(x_{t}^{\,s})-\nabla f_{i_t}(\tilde{x}^{\,s})+\tilde{\mu}^{\,s}$.  Adding and subtracting $\nabla f(x_{t}^{\,s})$ yields
\[
v_t^{\,s}= \underbrace{\nabla f_{i_t}(x_{t}^{\,s})-\nabla f(x_{t}^{\,s})}_{\xi_{1}}
-\underbrace{\bigl(\nabla f_{i_t}(\tilde{x}^{\,s})-\nabla f(\tilde{x}^{\,s})\bigr)}_{\xi_{2}}
+\nabla f(x_{t}^{\,s}).
\tag{11}
\]
Hence,
\[
\|v_t^{\,s}\|^{2}
\le 3\|\nabla f(x_{t}^{\,s})\|^{2}+3\|\xi_{1}\|^{2}+3\|\xi_{2}\|^{2}.
\tag{12}
\]
\textbf{[Step 2]} Taking expectation conditioned on $x_{t}^{\,s},\tilde{x}^{\,s}$ and using Assumption~\ref{ass:var},
\[
\EE\!\bigl[\|v_t^{\,s}\|^{2}\mid x_{t}^{\,s},\tilde{x}^{\,s}\bigr]
\le 3\|\nabla f(x_{t}^{\,s})\|^{2}+6\sigma^{2}.
\tag{13}
\]

\subsection*{Step 3: Combine (10) and (13)}
Plugging (13) into (10) gives
\[
\EE\!\bigl[f(x_{t+1}^{\,s})\mid x_{t}^{\,s}\bigr]
\le f(x_{t}^{\,s})-\eta\|\nabla f(x_{t}^{\,s})\|^{2}
+\frac{3L\eta^{2}}{2}\|\nabla f(x_{t}^{\,s})\|^{2}+3L\eta^{2}\sigma^{2}.
\tag{14}
\]
Rearrange:
\[
\EE\!\bigl[f(x_{t+1}^{\,s})\mid x_{t}^{\,s}\bigr]
\le f(x_{t}^{\,s})-\Bigl(\eta-\frac{3L\eta^{2}}{2}\Bigr)\|\nabla f(x_{t}^{\,s})\|^{2}
+3L\eta^{2}\sigma^{2}.
\tag{15}
\]

\subsection*{Step 4: Add the distance term}
From the definition of $\Phi_{t}^{\,s}$ (5) we also need to control $\|x_{t+1}^{\,s}-\tilde{x}^{\,s}\|^{2}$.
Using $x_{t+1}^{\,s}=x_{t}^{\,s}-\eta v_t^{\,s}$,
\[
\|x_{t+1}^{\,s}-\tilde{x}^{\,s}\|^{2}
= \|x_{t}^{\,s}-\tilde{x}^{\,s}\|^{2}
-2\eta\langle v_t^{\,s},x_{t}^{\,s}-\tilde{x}^{\,s}\rangle
+\eta^{2}\|v_t^{\,s}\|^{2}.
\tag{16}
\]
Take expectation conditioned on $x_{t}^{\,s},\tilde{x}^{\,s}$.  The cross term is handled by Cauchy–Schwarz and Young’s inequality:
\[
-2\eta\EE[\langle v_t^{\,s},x_{t}^{\,s}-\tilde{x}^{\,s}\rangle]
\le 2\eta\EE[\|v_t^{\,s}\|\|x_{t}^{\,s}-\tilde{x}^{\,s}\|]
\le \eta\EE[\|v_t^{\,s}\|^{2}]+ \eta\|x_{t}^{\,s}-\tilde{x}^{\,s}\|^{2}.
\tag{17}
\]
Thus,
\[
\EE\!\bigl[\|x_{t+1}^{\,s}-\tilde{x}^{\,s}\|^{2}\mid x_{t}^{\,s},\tilde{x}^{\,s}\bigr]
\le (1+\eta)\|x_{t}^{\,s}-\tilde{x}^{\,s}\|^{2}+2\eta^{2}\EE[\|v_t^{\,s}\|^{2}].
\tag{18}
\]
Insert the bound (13) for $\EE[\|v_t^{\,s}\|^{2}]$:
\[
\EE\!\bigl[\|x_{t+1}^{\,s}-\tilde{x}^{\,s}\|^{2}\mid x_{t}^{\,s},\tilde{x}^{\,s}\bigr]
\le (1+\eta)\|x_{t}^{\,s}-\tilde{x}^{\,s}\|^{2}
+6\eta^{2}\|\nabla f(x_{t}^{\,s})\|^{2}+12\eta^{2}\sigma^{2}.
\tag{19}
\]

\subsection*{Step 5: Lyapunov descent}
Multiply (19) by $c/2$ and add to (15).  Choosing $c:=\frac{12L\eta}{1-\eta}$ (which is positive for $\eta<1$) yields cancellation of the $\|\nabla f\|^{2}$ terms:
\[
\EE\!\bigl[\Phi_{t+1}^{\,s}\mid x_{t}^{\,s},\tilde{x}^{\,s}\bigr]
\le \Phi_{t}^{\,s}
-\underbrace{\Bigl(\eta-\frac{3L\eta^{2}}{2}-\frac{c\eta^{2}}{2}\cdot6\Bigr)}_{\displaystyle\ge\frac{\eta}{2}}\|\nabla f(x_{t}^{\,s})\|^{2}
+ \underbrace{\Bigl(3L\eta^{2}+ \frac{c}{2}\cdot12\eta^{2}\Bigr)}_{\displaystyle C\eta^{2}}\sigma^{2},
\tag{20}
\]
where the inequality $\eta-\frac{3L\eta^{2}}{2}-3c\eta^{2}\ge \eta/2$ holds whenever $\eta\le \frac{1}{3L}$ and $c$ is defined as above.  Denote $C:=3L+6c$.

\textbf{[Step 5]} Hence for all inner steps,
\[
\EE\!\bigl[\Phi_{t+1}^{\,s}\bigr]\le \EE\!\bigl[\Phi_{t}^{\,s}\bigr]-\frac{\eta}{2}\EE\!\bigl[\|\nabla f(x_{t}^{\,s})\|^{2}\bigr]+C\eta^{2}\sigma^{2}.
\tag{21}
\]

\subsection*{Step 6: Summation over an epoch}
Summing (21) for $t=0,\dots,m-1$ telescopes the Lyapunov function:
\[
\EE\!\bigl[\Phi_{m}^{\,s}\bigr]\le \EE\!\bigl[\Phi_{0}^{\,s}\bigr]
-\frac{\eta}{2}\sum_{t=0}^{m-1}\EE\!\bigl[\|\nabla f(x_{t}^{\,s})\|^{2}\bigr]
+ mC\eta^{2}\sigma^{2}.
\tag{22}
\]
Recall $\Phi_{0}^{\,s}= f(\tilde{x}^{\,s})+\frac{c}{2}\cdot0$ and $\Phi_{m}^{\,s}= f(\tilde{x}^{\,s+1})+\frac{c}{2}\| \tilde{x}^{\,s+1}-\tilde{x}^{\,s}\|^{2}\ge f(\tilde{x}^{\,s+1})$.
Thus,
\[
\frac{\eta}{2}\sum_{t=0}^{m-1}\EE\!\bigl[\|\nabla f(x_{t}^{\,s})\|^{2}\bigr]
\le f(\tilde{x}^{\,s})- \EE\!\bigl[f(\tilde{x}^{\,s+1})\bigr] + mC\eta^{2}\sigma^{2}.
\tag{23}
\]

\subsection*{Step 7: Sum over all epochs}
Let $S$ be the total number of epochs and $T=Sm$ the total stochastic gradient evaluations.  Summing (23) over $s=0,\dots,S-1$ gives
\[
\frac{\eta}{2}\sum_{s=0}^{S-1}\sum_{t=0}^{m-1}\EE\!\bigl[\|\nabla f(x_{t}^{\,s})\|^{2}\bigr]
\le f(\tilde{x}^{\,0})- \EE\!\bigl[f(\tilde{x}^{\,S})\bigr] + SmC\eta^{2}\sigma^{2}.
\tag{24}
\]
Since $f(\tilde{x}^{\,S})\ge f^{\star}$ and $Sm=T$,
\[
\frac{1}{T}\sum_{k=1}^{T}\EE\!\bigl[\|\nabla f(x_{k})\|^{2}\bigr]
\le \frac{2\bigl(f(\tilde{x}^{\,0})-f^{\star}\bigr)}{\eta T}+2C\eta\sigma^{2}.
\tag{25}
\]

\subsection*{Step 8: Choice of stepsize}
Set $\eta=1/(3L)$ and $m=2n$ (hence $c=12L\eta/(1-\eta)\le 6$).  With these values $C\le 9L$.
Plugging into (25) yields
\[
\EE\!\bigl[\|\nabla f(\hat{x})\|^{2}\bigr]
\le \frac{3L\bigl(f(\tilde{x}^{\,0})-f^{\star}\bigr)}{T}
+\frac{2\sigma^{2}}{n},
\tag{26}
\]
which is exactly the bound claimed in Theorem~\ref{thm:svrg}.  $\square$

\section*{Rate Derivation (With Constants)}
From (26) we see that the dominant term decays as $O(1/T)$.  To achieve $\EE[\|\nabla f(\hat{x})\|^{2}]\le\varepsilon$ it suffices to pick
\[
T\ge \frac{3L\Delta_f}{\varepsilon}\quad\text{and}\quad
n\ge \frac{2\sigma^{2}}{\varepsilon},
\]
or, more compactly,
\[
T = O\!\Bigl(\frac{L\Delta_f}{\varepsilon} + \frac{n\sigma^{2}}{\varepsilon}\Bigr).
\]

\section*{Special Cases}
\begin{itemize}
\item \textbf{Convex $L$‑smooth case.} The same proof works, and the bound (26) can be strengthened to an $O(1/T)$ guarantee on the objective suboptimality $f(\hat{x})-f^{\star}$.
\item \textbf{Infinite‑sum (online) setting.} When $n\to\infty$ and the variance bound $\sigma^{2}$ is finite, the term $2\sigma^{2}/n$ vanishes, yielding a pure $O(1/T)$ rate.
\item \textbf{Large inner loop $m\gg n$.} The constant $C$ grows linearly with $m$, which deteriorates the bound; thus $m=\Theta(n)$ is optimal.
\end{itemize}

\end{document}